\documentclass[10pt,a4paper]{article}

\usepackage{fontspec}
\usepackage[a4paper,margin=20mm,headheight=15pt]{geometry}
\usepackage{booktabs}
\usepackage{longtable}
\usepackage{array}
\usepackage{enumitem}
\usepackage{xcolor}
\usepackage{hyperref}
\usepackage{fancyvrb}
\usepackage{lscape}

\hypersetup{
  colorlinks=true,
  linkcolor=blue!55!black,
  urlcolor=blue!55!black,
  pdftitle={RICH Ablation Merge: Audited Change Report},
  pdfauthor={Codex and eladtan}
}

\setlist[itemize]{leftmargin=*,topsep=3pt,itemsep=2pt}
\setlist[enumerate]{leftmargin=*,topsep=3pt,itemsep=2pt}
\setlength{\LTpre}{4pt}
\setlength{\LTpost}{6pt}
\setlength{\parindent}{0pt}
\setlength{\parskip}{5pt}
\sloppy
\setmonofont{Source Code Pro}

\definecolor{passgreen}{RGB}{25,115,55}
\definecolor{warnorange}{RGB}{170,92,0}
\definecolor{codegray}{RGB}{248,248,248}
\newcommand{\pass}{\textcolor{passgreen}{\textbf{PASS}}}
\newcommand{\commit}[1]{\texttt{#1}}
\newcommand{\file}[1]{\nolinkurl{#1}}

\title{\textbf{RICH Ablation Merge}\\[3pt]
  Audited Change, Validation, and Diff Report}
\author{Prepared from the pushed Git objects by Codex}
\date{6 August 2026}

\begin{document}
\maketitle

\begin{center}
\begin{tabular}{@{}ll@{}}
\toprule
Repository & \file{/home/elads/RICH-ablation-integration} \\
Branch & \file{codex/merge-ablation-20260805} \\
Merge commit & \commit{6d6192cbb90bdfdaf750b8e33fefdd3bad8c657b} \\
First parent & \commit{7802210407bd416aef0f48781104328d07a6b926} \\
Ablation parent & \commit{8031e78e4b82f79a5189bf953fcfeab1086644f2} \\
STORM gitlink & \commit{184bd16f84f4c05424e675e9945b357592a21f5a} \\
THUNDER gitlink & \commit{5cd60ccec76eea829c25fe3136da236d23fe0e41} \\
\bottomrule
\end{tabular}
\end{center}

\tableofcontents
\clearpage

\section{Scope and Method}

This report describes the effective first-parent change set of merge commit
\commit{6d6192cb}.  That commit merges the ablation parent
\commit{8031e78e} into the integration baseline \commit{78022104}.  Therefore,
the complete diff contains both imported ablation work and repairs made while
integrating, compiling, and validating it.  The report distinguishes those
classes where the Git history and work log permit it; it does not claim that
all imported code was authored during the integration session.

The audit used immutable Git objects rather than the current working-tree
contents:
\begin{itemize}
  \item parent diff: \commit{78022104..6d6192cb};
  \item expanded STORM diff: \commit{9f2b9fd..184bd16};
  \item expanded THUNDER diff: \commit{8e82079..5cd60cc};
  \item regression artifacts under \file{regression_results}; and
  \item Slurm accounting for the build and THUNDER-unit jobs.
\end{itemize}

The final appendices contain the complete file-status diff, complete numstat,
and complete textual unified patches.  The two changed JPEG files are recorded
by Git as binary differences; their raw binary bytes are intentionally not
printed as text.

\section{Executive Summary}

The merge adds the serial and distributed fast-multipole gravity stack,
substantial radiation post-processing and polarization infrastructure,
DDMC moving-interface and cross-rank hardening, spherical finite-volume
geometry support, expanded regression coverage, and many production/run
configurations.  Integration-specific work then made compiler/MPI/VTK
selection portable, fixed asymmetric DDMC ghost exchange, hardened THUNDER's
Slurm execution model, completed the STORM port, and removed the unwanted
reverse-estimator implementation.

The first-parent diff contains:
\begin{center}
\begin{tabular}{@{}rrrrrr@{}}
\toprule
Files & Added & Modified & Deleted & Renamed & Text delta \\
\midrule
238 & 141 & 81 & 14 & 2 & $+51{,}247/-10{,}813$ \\
\bottomrule
\end{tabular}
\end{center}

The final required regression set, after the explicitly removed tests and the
seven explicitly excluded large/manual cases, is \textbf{51/51 passing with
GNU} and \textbf{51/51 passing with Intel OneAPI plus OpenMPI}.  No merge into
the repository's main/master branch was performed.

\section{Change Inventory by Area}

\begin{longtable}{@{}>{\raggedright\arraybackslash}p{55mm}rrrr@{}}
\toprule
Area & Files & Insertions & Deletions & Binary \\
\midrule
\endhead
Build, repository, and miscellaneous & 6 & 127 & 111 & 0 \\
Analysis and migration tools & 5 & 2,285 & 0 & 0 \\
Documentation & 18 & 1,225 & 73 & 0 \\
Regression infrastructure and cases & 57 & 9,448 & 143 & 0 \\
Run configurations and reference outputs & 30 & 3,151 & 5,905 & 2 \\
Gravity and FMM & 54 & 21,278 & 24 & 0 \\
Radiation and post-processing & 50 & 11,614 & 4,357 & 0 \\
Hydro, geometry, and mesh & 17 & 2,118 & 199 & 0 \\
STORM submodule gitlink & 1 & 1 & 1 & 0 \\
\midrule
\textbf{Total} & \textbf{238} & \textbf{51,247} & \textbf{10,813} & \textbf{2} \\
\bottomrule
\end{longtable}

The THUNDER gitlink is counted in the regression row because its path is under
\file{regression_tests}; the STORM gitlink is shown separately.

\section{Build, Compiler, MPI, and VTK Portability}

\subsection{Compiler and MPI selection}

\file{config/set_compilers.cmake} was changed so that compiler/module changes
do not inherit stale FindMPI state from a reused build directory.  The final
logic:
\begin{itemize}
  \item resolves \file{mpiexec}, \file{mpicc}, \file{mpicxx}, and
        \file{mpifort} from the active \file{PATH};
  \item clears derived MPI cache variables before re-probing;
  \item configures OpenMPI builds through OpenMPI wrappers, including when the
        underlying compiler is Intel OneAPI;
  \item validates the wrapper backend instead of assuming it from the compiler
        executable's name; and
  \item avoids hard-coded cluster paths, host names, or module names in source.
\end{itemize}

This permits both GCC/OpenMPI and Intel-OneAPI/OpenMPI configurations without
binding the source to the validation machine.

\subsection{VTK/MPI compatibility}

\file{config/find_VTK.cmake} now seeds transitive VTK FindMPI calls from the
active wrappers and clears stale cached results.  For OpenMPI builds it checks
the actual shared-library dependency of \file{VTK::ParallelMPI}, because a
compile-only \file{vtkMPICommunicatorOpaqueComm} probe is not a reliable MPI
vendor test.  Binary inspection prefers CMake's toolchain-provided
\file{CMAKE_READELF}, falls back to \file{readelf}/\file{llvm-readelf}, and
degrades to a warning where no compatible inspector exists.  This keeps the
fix portable beyond ELF-only environments.

\section{Gravity and Fast Multipole Method}

The merge introduces a complete FMM implementation under
\file{source/3D/gravity/fmm} and connects it to the production three-dimensional
gravity interface.

\begin{itemize}
  \item \textbf{Serial solver:} tree construction, multipole/local expansion
        layout, Laplace solid harmonics, P2M/M2M/M2L/L2L/L2P passes, dual-tree
        traversal, operator caching, root geometry, and direct-reference
        validation.
  \item \textbf{Distributed solver:} global dyadic lattice, local patch and
        patch-forest representations, process trees, LET planning, peer packet
        exchange, descriptor gathering, traversal partitioning, persistent
        topology state, and collective validation of distributed options.
  \item \textbf{Production integration:}
        \file{FastMultipoleAcceleration3D} and updates to distributed gravity
        calculators expose FMM configuration and diagnostics while preserving
        the existing gravity entry points.
  \item \textbf{Diagnostics:} timing, memory, root mass, leaf occupancy,
        interaction counts, plan reuse/rebuild, cache behavior, rejection
        reasons, and direct-sum error measurements.
  \item \textbf{Analysis tooling:} new scripts inspect patch traces, solve
        traces, topology nonuniformity, and transformation/migration behavior.
\end{itemize}

Regression cases cover serial accuracy, MPI construction guards, distributed
accuracy with empty ranks, process-pair coverage, peer-plan rebuilds, operator
cache behavior, and the FMM/quadrupole benchmark.  The two very large
lane/benchmark cases remain outside the required capped matrix where specified
by the user.

\section{Radiation, DDMC, and Post-processing}

\subsection{Forward post-processing and polarization}

The radiation post-processing code is reorganized under
\file{source/3D/radiation/postprocess}.  New components implement adaptive
statistics, flux-source and grey calculations, photosphere calculations,
command-line/configuration handling, runtime state, result writing, and the
post-processing application.  Existing helper source/header files were moved
into that directory with 98\% and 97\% rename similarity.

Observer, polarization, and Lorentz-transformation code was extended, including
new polarization statistics and updates to spherical-observer handling.
Production and run-level post-processing controls were updated accordingly.

\subsection{Reverse estimator removed}

The reverse estimator was deliberately removed rather than retained as an
optional mode.  Twelve implementation/header files were deleted:
\begin{itemize}
  \item \file{ReverseAdjointTransport3D.cpp/.hpp},
        \file{ReverseDDMC.cpp/.hpp}, and \file{ReverseDoppler.cpp/.hpp};
  \item \file{ReverseEstimatorConfig.hpp}, \file{ReversePacket.hpp};
  \item \file{ReverseObserverTallies.cpp/.hpp}; and
  \item \file{ReversePolarizationMueller.cpp/.hpp}.
\end{itemize}
Associated CMake registrations, command-line/configuration controls, object
libraries, tests, and submission options were removed.  A final source-tree
search found no reverse-estimator reference.

\subsection{DDMC interface and distributed exchange}

The merge adds the Wollaeger moving-interface kernel and diagnostics, plus
documentation of incidence, admission/reflection, moving factors, fallback,
splitting, and reciprocity.  Cross-rank DDMC validation was retained through
\file{ddmc_mpi_zero_cell}.

The integration repair in STORM's
\file{radiation/ddmc/DDMCGhostExchange.hpp} addresses asymmetric domain
decompositions.  Storage is sized for local plus ghost cells before indexed
access; owner aggregation and ghost clearing use their own valid lengths
rather than assuming symmetric ownership.  A focused GNU MPI regression
confirmed the asymmetric reduction path.

The user later removed \file{spherical_density_hardening_ddmc_push} and
\file{ddmc_moving_interface_ab}.  Their case directories, checker functions,
standalone checker/wrapper, and documentation/catalog associations are absent
from the final tree.  The already-running spherical production jobs were
canceled because the test was no longer part of the requested product.

\section{STORM Submodule}

The parent gitlink advances STORM from \commit{9f2b9fd} to
\commit{184bd16}.  The expanded range contains two commits and changes five
files by $+1{,}866/-42$:
\begin{itemize}
  \item port ablation radiation controls into \file{RadiationIMC.hpp},
        \file{RadiationIMCParameters.hpp}, \file{DDMCTypes.hpp}, and the DDMC
        exchange layer;
  \item make pre-project MPI compiler selection work through the active
        wrappers in \file{CMakeLists.txt};
  \item add compatibility dispatch for equation-of-state \file{dT2e}
        signatures; and
  \item fix the asymmetric DDMC ghost/owner exchange described above.
\end{itemize}

The final submodule commit is pushed on
\file{codex/ablation-port-20260805}.

\section{THUNDER Submodule and Regression Infrastructure}

THUNDER advances from \commit{8e82079} to \commit{5cd60cc}; two files change
by $+161/-6$.
\begin{itemize}
  \item arithmetic-looking metadata is resolved safely before integer
        conversion, eliminating the crash caused by a literal shell expression
        in task-count metadata;
  \item compiler/profile/test-specific build directories prevent concurrent or
        sequential profiles from reusing incompatible CMake state;
  \item Slurm time-limit parsing accepts scheduler formats and respects
        \file{SLURM_TIME_LIMIT};
  \item fast jobs are not lost in the submit/poll transition; and
  \item unit coverage was expanded for these behaviors.
\end{itemize}

The THUNDER unit suite passed 20/20 on Slurm job \commit{10104960}.  The final
submodule commit is pushed on
\file{codex/thunder-portability-20260806}.

The parent regression layer also gains dependency-aware aggregate checks,
expanded checker functions, result/report tooling, FMM and DDMC cases,
radiation-direction sampling, and Densmore interface diagnostics.  The removed
\file{spherical_collapse_hires} case is recorded as a deletion in the complete
diff.

\section{Hydrodynamics, Geometry, and Mesh Work}

New spherical finite-volume support includes spherical-shell geometry and
projection classes, a spherical-background linear-Gauss implementation, and
updates to the existing spherical reconstruction path.  AMR, neighbor,
tessellation I/O, mesh-generation, and universal-error paths were updated to
support the merged gravity/radiation behavior.  Radiation and hydrodynamic
simulation steps were updated where the merged interfaces changed.

\section{Run Configurations, Documentation, and Utilities}

Run-level changes add or revise Compton, Marshak, gold heat-wave, moving-slab,
and TDE post-processing configurations.  Several diagnostic and plotting
utilities were added for polarization, spectra, Densmore interfaces, and FMM
behavior.  Two JPEG comparison artifacts changed and are represented as
binary differences in the patch.

Documentation was expanded for FMM architecture, MPI FMM and patch forests,
DDMC/Wollaeger behavior, gravity user controls, regression execution, and
merge hazards.  Wiki copies of regression documentation were updated with the
same final behavior.

\section{Validation and Self-check}

\subsection{Compiler matrices}

All builds and executions used the \file{bigrun} partition and no required test
used more than eight nodes.

\begin{longtable}{@{}p{40mm}p{28mm}p{77mm}@{}}
\toprule
Matrix & Result & Evidence/reconciliation \\
\midrule
\endhead
GNU Release, serial then MPI & \pass & Main artifact: 52/52.  The subsequently
removed \file{ddmc_moving_interface_ab} had passed, leaving the final required
set at 51/51. \\
Intel OneAPI 2024.2.1 with OpenMPI & \pass & Main artifact: 48/53.  Three STORM
wrapper retries and one FMM benchmark retry passed.  Of the original five
non-passes, \file{spherical_density_hardening_ddmc_push} was removed; the
already-passing \file{ddmc_moving_interface_ab} was also removed.  Final
required set: 51/51. \\
THUNDER unit suite & \pass & 20/20, Slurm job 10104960. \\
Focused asymmetric DDMC MPI case & \pass & GNU retry artifact passed after the
ghost-exchange repair. \\
\bottomrule
\end{longtable}

GNU and Intel build jobs \commit{10104979} and \commit{10104984} completed with
exit code 0 on \file{bigrun}.  Both formerly suspect linear interpolation cases,
\file{cartesian_gauss_linear} and \file{spherical_gauss_linear}, passed with
both compiler stacks.

Exactly seven large/manual cases were excluded from the required matrix:
\file{lane_self_gravity_fmm}, \file{spherical_collapse_hires},
\file{hohlraum_parallel}, \file{lane_self_gravity}, \file{moving_slab},
\file{fmm_mpi_scaling_benchmark}, and \file{moving_slab_mc_32}.

\subsection{Static and provenance checks}

The final self-check found:
\begin{itemize}
  \item no whitespace errors in the parent, STORM, or THUNDER diffs;
  \item no unresolved Git entries or merge-conflict markers;
  \item no source references to the reverse estimator or the two explicitly
        removed tests;
  \item no validation-machine paths, host names, or module commands added by
        the compiler/STORM/THUNDER portability fixes;
  \item pushed, fetchable submodule commits at the exact gitlinks recorded by
        the parent; and
  \item historical validation artifacts retained separately from source.
\end{itemize}

Four unrelated pre-existing working-tree deletions
(\file{asan_mpi.supp}, \file{asan_wrap.sh},
\file{imc_compton_implementation_plan.tex}, and \file{mainpage.h}) were not
included in the merge commit.

\section{Conclusions and Known Limits}

The final integration branch is internally consistent, portable across the
validated GNU and Intel/OpenMPI configurations, and backed by a complete
passing required regression matrix.  The report does not claim runtime
validation of the seven explicitly excluded cases.  It also does not treat
the removed spherical/DDMC A/B tests as requirements.  No merge into a mainline
branch or merge request was created.

\clearpage
\appendix
\section{Complete File-status Diff}

The following is the exact \file{git diff --name-status --find-renames} output
for \commit{78022104..6d6192cb}.

\begin{landscape}
\VerbatimInput[
  fontsize=\tiny,
  frame=single,
  numbers=left,
  numbersep=3pt
]{merge_ablation_file_status.txt}
\end{landscape}

\section{Complete File Numstat}

The columns are inserted lines, deleted lines, and path.  A dash denotes a
binary file.

\begin{landscape}
\VerbatimInput[
  fontsize=\tiny,
  frame=single,
  numbers=left,
  numbersep=3pt
]{merge_ablation_file_numstat.txt}
\end{landscape}

\section{Complete Unified Patches}

These patches are the exact textual Git diffs used for this audit.  The parent
patch is first-parent based; the STORM and THUNDER patches expand the submodule
gitlink changes so their actual source modifications appear at the end of this
document.

\subsection{Parent repository: \commit{78022104..6d6192cb}}
\begin{landscape}
\VerbatimInput[
  fontsize=\tiny,
  frame=single,
  numbers=left,
  numbersep=3pt,
  tabsize=2
]{merge_ablation_parent.patch}
\end{landscape}

\subsection{STORM: \commit{9f2b9fd..184bd16}}
\begin{landscape}
\VerbatimInput[
  fontsize=\tiny,
  frame=single,
  numbers=left,
  numbersep=3pt,
  tabsize=2
]{merge_ablation_storm.patch}
\end{landscape}

\subsection{THUNDER: \commit{8e82079..5cd60cc}}
\begin{landscape}
\VerbatimInput[
  fontsize=\tiny,
  frame=single,
  numbers=left,
  numbersep=3pt,
  tabsize=2
]{merge_ablation_thunder.patch}
\end{landscape}

\end{document}
